\documentclass[12pt]{article}

\usepackage[a4paper, total={6in, 8in}]{geometry}
\usepackage{textcomp}

\begin{document}
\setlength{\parindent}{0pt}

\def \t {\textrightarrow}
\def \v {\vspace{1em}}
\def \bi {\begin{itemize}}
\def \ei {\end{itemize}}
\def \s[#1] {\section*{#1}}
\def \ss[#1] {\subsection*{#1}}
\def \sss[#1] {\subsubsection*{#1}}

\s[Umberto Saba - Amai]
Sintetizza la sua visione della vita e della poesia \t la poesia delle cose semplici, una poesia etica e fatta di oggetti, di immagini quotidiane senza orpelli d'annunziani \\
Il testo è il testamento poetico dell'autore \\
Il testo ha una caratteristica importante, ribadita da anafora verso 1, 5 e 9 \\
Come se la vita fosse un gioco, che si realizza in un continuo contrasto \t opera tratta dal canzoniere \\
Il suo romanzo Ernesto \t che ha tema dell'amore omosessuale e rapporto vissuto tra autore e ragazzo

\v

Amai \t che fa al fondo, poeta palombaro \t l'abisso, cercare la verita che giace al fondo \\
Il poeta che ricerca la poesia onesta esprime la personalita di Umberto Saba \\
Compaiono i fantasmi che l'analista ??????? \t il sogno con Freud non è premonitore, ma una manifestazione di un irrisolto \t ma non ha valore premonitore \\
"Quasi un sogno bliato" \t un sogno che non si ritrova nella realta \\
Primo verso fondamentale \\
Amo te \t te siamo noi, il lettore \t in questo si legge il rapporto del poeta \t comunica semplicemente un rapporto con il lettore, dice di avere qualcosa da dire al lettore \\
È il testamento dell'autore, che si manifesta attraverso l'idea di questa carta e del gioco che è la vita \t anche Montale lascia testamento, prassi inaugurata da Petrarca (lettera ai posteri) \\
Poliptoto tra amai e amo \t la vita come gioco, verso chiave da cui partire è "la verità che giace al fondo" = tutto il messaggio di saba, il poeta è colui che scava negli abissi ed è il palombaro di Ungaretti \\
La poesia precede la psicanalisi \t Grossman dice psicanalitico prima della psicanalisi

\v

Il titolo del canzoniere presenta un radicamento nel passato e uno slancio nel futuro \\
La transizione petrarchesca e questo volto \t gli antichi avevano gia detto tutto, e petrarca che vive il senso del tempo, un tempo interiore che Bergson descrive nella differenza tra il tempo interiore e quello della realta esterna \\
Il tempo interiore non ha convenzioni se non con noi stessi

\v

Questo incentra su giudizi e pregiudizi sull'omosessualità \t saba ha inclinazione omosessuale \t Vittorio Lingiardi si occupa di questo fenomeno sugli intellettuali \\
Pasolini è stato uno dei maggiori critici di saba \t è come un io giudicante che ha vissuto \\
Saba è attuale per la sua quotidianità e semplicita apparente \t dignita del quotidiano che non risparmia aspetti + tragici: il padre è un assassino \\
Il legame con Beppa non è sanguigno ma passa attraverso il latte \t è anche un nutrimento dell'anima \\
L'ernesto propone tematica del romanzo, in cui non c'è morbosita \t non è presente, e l'omosessualità viene vissuta con naturalezza \\
Saba muore nel 57 a Gorizia

\s[Crepuscolarismo]
Inizio 900 c'è crisi dell'io, e il nuovo secolo presenta innovazioni che lasciano certi agi e comfort \t treno nella stazione di una mattina d'autunno es. \\
Nascono due movimenti distinti: il futurismo e il crepuscolarismo \\
Crepuscolarismo \t si sviluppa dal 1902, e viene definito da Borgese (che parla anche di d'annunzio) \t movimento evoca il momento della giornata che precede la sera, che genera senzo di malinconia e ripiegamento interiore \\
Questi autori hanno sentimento comune di volersi allontanare da d'annunzio, e di recuperare una poesia di "piccole cose di pessimo gusto" \t che genera oggettistica \\
L'oggettistica di d'annunzio è quella degli orpelli, quella di saba sono le cose semplici, quella di pascoli sono uccelli \t del crepuscolarismo sono i salotti dell'800, loreto impagliato, busto di alfieri, merletti, e figure femminili che sono la Muti del piacere e altre sensuali \\
Come Felicita, compagna di Guido Cozzano, che hanno gli occhi color stoviglia (azzurro spento) \\
L'oggettistica crepuscoalre comprende anche paesaggi come giardini chiuse, fontane asciute, organi che suonano, corsie d'ospedale = anti d'annunzio, 1910 definizione di crepuscolarismo \\
Ma in tutto questo mondo, il poeta non è + il vate \t allora cos'è?

\v

Sergio Corazzini \t desolazione di un povero poeta senimentale \t il poeta è il bambino che piange \t che è diverso dal fanciullo di pascoli \\
Aldo Palazzeschi rappresenta invece il gancio tra crepuscolari e futuristi \t legati tra loro in sinergia, e vivono il dramma della confusione di inizio 900 \\
Il volto malinconico del crepuscolarismo, Palazzeschi sposa anche visione futurista di "incendiare il mondo" \t si fossi foco, cecco angolieri \\
La mania piromane del futurismo di palazzeschi \t prima raccolta è "I cavalli bianchi" 1902 \t qua le immagini apparentemente cantilenanti, un mondo ovattato che nasconde dalla societa giudicante \\
Poi pubblica anche "La laterna" del 1907, 1909 manifesto di Marinetti sul Figaro, quindi c'è coincidenza \\
Si distrugge il chiaro di luna, la sintassi, le biblioteche, il passatismo, il manzonismo \t e cosa costruiamo? \\
Edizione del 09 e del 11 \t nome di Levi, su cui ha lavorato, che dimostra che Palazzeschi è di fatto un decadente \t ne crepuscolare ne futurista, aleggiano quindi Iam Rotterdamn e Metterlinch

\v

Il futurismo distrugge ma non costruisce \t la piromania significa ardere tutto, mettere tutto in una fiamma \\
La scapgigliatura e i poeti maledetti \t questa espressione di protesta ci pone di fronte alla volonta di costuire una nuova era \t è l ora dell'anticristo \\
Qua invece non c'è spirito costruttivo \t un incendio che brucia lascia una cenere e non contempla araba fenice \t questo non vale pero con il crepuscolarismo

\v

Il chiaro di luna fa parte degli stereotipi romantici \t è il chiaro di luna dell'addio ai monti, o di Debussy \t mentre qua la luce non puo piu essere accettata \\
C'è il teatro futurista, le serate futuriste, etc. \t è un fenomeno di costume \t ma cosa lascia? \t lascia in eredita un nuovo equilibrio tra spazi bianchi e parola scritta = tratto distintivo di Ungaretti \\
Ungaretti porta nella lirica questo nuovo contesto di testualita \t una nuova armonia tra spazi bianchi e parola scritta \\
Le parole non rispettano la punteggiatura \t c'è liberta compositiva che crea non senso che è il senso del futurista \\
Le parole in libertà \t libertà che è testuale, contiene in realta la svolta ungerettiana della poesia pura, dove la parola si è epurata delle scorie d'annunziane, e come parola pura indaga scandaglio interiore \\
C'è senso di sgomento di una società messa al rogo \t Marinetti pubblica manifesto che coinvolge l'arte e una mentalita dissacrante e piromane che si configura nel futurismo \\
Papini afferma che la guerra è l'unica igene del mondo \t serve umanita nuova, e distruggere quella vecchia violentemente

\ss[Palazzeschi]
Compone fino a tardi e abbraccia tutto il 900 \t ha anche lato futurista, e il suo aspetto futurista raggiunge l'apice nella pubblicazione del romano "Il codice di perllà" = romanzo fatto di tutte onomatopee \\
Non c'è senso logico, ma sovraccarico di onomatopee che esprimono il senso di Perelà, che è un uomo di fumo \\
L'assenza di una sintassi, di un nesso tematico, di una declinazione continua dell'onomatopea, fa capire come la realta è del tutto disintegrata \\
Sia nel "L'incendiaro", sia in questo romanzo, si coglie quanto il non senso (secondo le voci di Mengardi, Gioanola e Levì) e si definisce il trio d'union con un allegria dissacrante \\
Abbiamo l'immagine del riso amaro in Pirandello \t l'anziana signora, e l'avvertimento del contrario \t il sentimento del contrario \\
Subentra poi uno spirito + empatizzante, e il sentimento del contrario \t l'empatizzare con l'interlocutore descrive il binomio che si coglie con l'espressione riso amaro (sinestesia) \\
Il riso amaro si ricollega a figura del clown, che fa ridere ma cela tristezza \t in palazzeschi la tristezza forse per la sua omosessualità ?

\v

Futurismo in arte è un avanguardia, ma in letteratura no \t le avanguardie in letteratura sono quelle degli anni 60
\end{document}
